% Options for packages loaded elsewhere
\PassOptionsToPackage{unicode}{hyperref}
\PassOptionsToPackage{hyphens}{url}
%
\documentclass[
]{article}
\usepackage{amsmath,amssymb}
\usepackage{lmodern}
\usepackage{iftex}
\ifPDFTeX
  \usepackage[T1]{fontenc}
  \usepackage[utf8]{inputenc}
  \usepackage{textcomp} % provide euro and other symbols
\else % if luatex or xetex
  \usepackage{unicode-math}
  \defaultfontfeatures{Scale=MatchLowercase}
  \defaultfontfeatures[\rmfamily]{Ligatures=TeX,Scale=1}
\fi
% Use upquote if available, for straight quotes in verbatim environments
\IfFileExists{upquote.sty}{\usepackage{upquote}}{}
\IfFileExists{microtype.sty}{% use microtype if available
  \usepackage[]{microtype}
  \UseMicrotypeSet[protrusion]{basicmath} % disable protrusion for tt fonts
}{}
\makeatletter
\@ifundefined{KOMAClassName}{% if non-KOMA class
  \IfFileExists{parskip.sty}{%
    \usepackage{parskip}
  }{% else
    \setlength{\parindent}{0pt}
    \setlength{\parskip}{6pt plus 2pt minus 1pt}}
}{% if KOMA class
  \KOMAoptions{parskip=half}}
\makeatother
\usepackage{xcolor}
\IfFileExists{xurl.sty}{\usepackage{xurl}}{} % add URL line breaks if available
\IfFileExists{bookmark.sty}{\usepackage{bookmark}}{\usepackage{hyperref}}
\hypersetup{
  pdftitle={Panel-Based Evaluation of Interdisciplinary Researchers: A Protocol for National Agencies},
  pdfauthor={A. Rivero},
  hidelinks,
  pdfcreator={LaTeX via pandoc}}
\urlstyle{same} % disable monospaced font for URLs
\usepackage{longtable,booktabs,array}
\usepackage{calc} % for calculating minipage widths
% Correct order of tables after \paragraph or \subparagraph
\usepackage{etoolbox}
\makeatletter
\patchcmd\longtable{\par}{\if@noskipsec\mbox{}\fi\par}{}{}
\makeatother
% Allow footnotes in longtable head/foot
\IfFileExists{footnotehyper.sty}{\usepackage{footnotehyper}}{\usepackage{footnote}}
\makesavenoteenv{longtable}
\setlength{\emergencystretch}{3em} % prevent overfull lines
\providecommand{\tightlist}{%
  \setlength{\itemsep}{0pt}\setlength{\parskip}{0pt}}
\setcounter{secnumdepth}{-\maxdimen} % remove section numbering
\ifLuaTeX
  \usepackage{selnolig}  % disable illegal ligatures
\fi

\title{Panel-Based Evaluation of Interdisciplinary Researchers: A
Protocol for National Agencies}
\author{A. Rivero}
\date{2026}

\begin{document}
\maketitle
\begin{abstract}
National evaluation agencies face a structural challenge when assessing
researchers classified as ``Interdisciplinary'': standard
single-discipline panels cannot fairly evaluate portfolios that span
multiple fields. We propose an evaluation protocol based on a
three-component indicator panel --- disciplinary diversity, network
coherence, and cross-field effect --- that characterizes the type and
degree of boundary-crossing in a researcher's work. Using mock
researcher profiles, we demonstrate how the panel informs committee
composition and distinguishes genuine integrators from polymaths and
misclassified specialists. We identify six failure modes in
interdisciplinary evaluation and propose safeguards for each. The
protocol is designed to complement, not replace, expert judgment.
\end{abstract}

\hypertarget{introduction}{%
\section{Introduction}\label{introduction}}

National evaluation agencies periodically assess researchers for career
advancement, funding eligibility, and institutional accreditation. These
assessments typically rely on disciplinary panels: a committee of
experts in the researcher's field evaluates their publication record,
citations, and broader contributions. This procedure works well when the
researcher's portfolio falls within a single recognized discipline.

A problem arises when a researcher is classified in the science group
``Interdisciplinary'' --- a category that exists in several national
systems (e.g., Spain's ANECA, Italy's ANVUR) for researchers whose work
does not fit any single disciplinary panel. Standard evaluation
procedures assign reviewers from a single discipline, creating a
structural mismatch: the committee lacks expertise in part of the
researcher's portfolio, or worse, applies disciplinary norms
(publication venues, citation rates, methodological standards) that are
inappropriate for cross-disciplinary work.

This note proposes an evaluation protocol that uses a three-component
indicator panel to characterize interdisciplinary researchers and inform
committee composition. The panel, developed in a companion review
(Rivero, 2026), combines disciplinary diversity (\(\Delta\)), network
coherence (\(S\)), and cross-field effect (\(E\)). We demonstrate the
protocol on mock researcher profiles, identify failure modes, and
propose safeguards.

Operationally, we treat ``interdisciplinary'' as an integration claim,
not just a breadth claim: high input diversity can represent either
genuine integration or disconnected multidisciplinarity, and the
protocol is built to distinguish these cases explicitly.

\hypertarget{the-indicator-panel}{%
\section{The Indicator Panel}\label{the-indicator-panel}}

We briefly summarize the three panel components; formal definitions and
validation appear in the companion review.

\textbf{Diversity} (\(\Delta\)): The Rao-Stirling index measures the
disciplinary spread of a researcher's cited references, incorporating
variety, balance, and disparity. High \(\Delta\) indicates that the
researcher draws on a broad range of fields.

\textbf{Coherence} (\(S\)): The mean linkage strength measures the
average pairwise bibliographic coupling between a researcher's
publications. High \(S\) indicates that publications share references
across category boundaries --- the researcher's diverse inputs are woven
into a unified research program.

\textbf{Cross-field effect} (\(E\)): The fraction of citations received
from outside the researcher's primary category. High \(E\) indicates
that the researcher's work produces impact across disciplinary
boundaries.

The key insight is that no single component suffices. Diversity alone
cannot distinguish a genuine integrator (\(\Delta\) high, \(S\) high,
\(E\) high) from a polymath who publishes in multiple unrelated fields
(\(\Delta\) high, \(S\) near zero, \(E\) low). The full triple is
needed.

\hypertarget{evaluation-protocol}{%
\section{Evaluation Protocol}\label{evaluation-protocol}}

\hypertarget{step-1-compute-the-panel}{%
\subsection{Step 1: Compute the Panel}\label{step-1-compute-the-panel}}

For each researcher classified as ``Interdisciplinary,'' the agency
computes (\(\Delta\), \(S\), \(E\)) from publication and citation data.
This requires:

\begin{itemize}
\tightlist
\item
  A disciplinary classification of cited references (e.g., Web of
  Science subject categories or Scopus ASJC codes).
\item
  A pairwise similarity matrix between categories.
\item
  Citation data for the cross-field effect.
\end{itemize}

\hypertarget{step-2-classify-the-profile}{%
\subsection{Step 2: Classify the
Profile}\label{step-2-classify-the-profile}}

The panel values are compared against classification thresholds:

\begin{longtable}[]{@{}llll@{}}
\toprule
Classification & \(\Delta\) & \(S\) & \(E\) \\
\midrule
\endhead
Genuine integrator & \(\geq 0.40\) & \(\geq 0.30\) & \(\geq 0.30\) \\
Polymath (non-integrative) & \(\geq 0.40\) & \(< 0.15\) & \(< 0.15\) \\
Specialist (reclassify) & \(< 0.35\) & any & any \\
Provisional (CI overlap) & boundary-overlap & boundary-overlap &
boundary-overlap \\
Ambiguous (requires full panel review) & all other combinations & & \\
\bottomrule
\end{longtable}

These thresholds are illustrative and should be calibrated against
empirical distributions before operational deployment. When confidence
intervals overlap multiple rows, the case should be treated as
provisional and sent to full-panel qualitative review.

\hypertarget{step-3-compose-the-committee}{%
\subsection{Step 3: Compose the
Committee}\label{step-3-compose-the-committee}}

\textbf{Procedure.} Given researcher \(r\) with category proportion
vector \(p_r\) (the fraction of references in each category):

\begin{enumerate}
\def\labelenumi{\arabic{enumi}.}
\tightlist
\item
  Identify the primary categories: \(K_r = \{i : p_{r,i} \geq \tau\}\),
  where \(\tau = 0.15\).
\item
  For each category \(i \in K_r\), include at least one evaluator from
  discipline \(i\).
\item
  Include at least one evaluator with demonstrated cross-disciplinary
  expertise (\(\Delta > 0.40\)).
\item
  Committee size: \(|K_r| + 1\) (disciplines represented plus a
  cross-disciplinary chair).
\item
  If profile classification is provisional (CI overlap), add one methods
  evaluator tasked with uncertainty and robustness review.
\end{enumerate}

The committee composition is thus data-driven: it reflects the
researcher's actual disciplinary profile rather than an arbitrary
assignment.

\hypertarget{demonstration-on-mock-profiles}{%
\subsection{Demonstration on Mock
Profiles}\label{demonstration-on-mock-profiles}}

We illustrate the protocol with three mock researchers, all classified
as ``Interdisciplinary'' by the agency.

\begin{longtable}[]{@{}lllll@{}}
\toprule
Researcher & \(\Delta\) & \(S\) & \(E\) & Classification \\
\midrule
\endhead
Dr.~A (integrator) & 0.558 & 0.733 & 0.600 & Genuine integrator \\
Dr.~B (polymath) & 0.562 & 0.000 & 0.063 & Polymath \\
Dr.~D (specialist) & 0.288 & 0.881 & 0.211 & Specialist (reclassify) \\
\bottomrule
\end{longtable}

\textbf{Dr.~A} publishes across condensed matter physics, physical
chemistry, and molecular biology. Her publications share references
across category boundaries (\(S = 0.733\)), and her work is cited across
disciplines (\(E = 0.600\)). The panel correctly identifies her as a
genuine integrator. Note that Dr.~A and Dr.~B have nearly identical
diversity scores (0.558 vs. 0.562), yet receive opposite classifications
--- diversity alone is insufficient. The committee should include
evaluators from her three primary fields plus a cross-disciplinary chair
(4 members).

\textbf{Dr.~B} has published in five different fields, but each
publication is a single-field contribution with no shared references
across domains (\(S = 0.000\)). His work is cited almost exclusively
within each publication's own field (\(E = 0.063\)). The panel
identifies him as a polymath. He should be reclassified to his strongest
field or evaluated separately in each field --- the
``Interdisciplinary'' label is misleading.

\textbf{Dr.~D} concentrates in condensed matter physics and materials
science (\(\Delta = 0.288\)), neighboring fields with high mutual
similarity. Her coherence is high because all publications draw on the
same knowledge base. The panel identifies her as a misclassified
specialist who should be evaluated by a standard disciplinary panel for
condensed matter physics.

\hypertarget{failure-modes-and-safeguards}{%
\section{Failure Modes and
Safeguards}\label{failure-modes-and-safeguards}}

We identify six failure modes that can arise in interdisciplinary
evaluation, even when using the panel.

\hypertarget{f1-breadth-without-depth-reward}{%
\subsection{F1: Breadth-without-depth
reward}\label{f1-breadth-without-depth-reward}}

\textbf{Risk:} A researcher with high \(\Delta\) is rewarded for breadth
regardless of coherence or impact.

\textbf{Safeguard:} Require \(S \geq 0.30\) and \(E \geq 0.30\) for
``integrator'' classification. High diversity alone does not qualify.

\hypertarget{f2-non-standard-publication-penalty}{%
\subsection{F2: Non-standard publication
penalty}\label{f2-non-standard-publication-penalty}}

\textbf{Risk:} Interdisciplinary journals typically have lower impact
factors than top disciplinary journals. Researchers publishing in such
venues are penalized in impact-factor-based assessments.

\textbf{Safeguard:} Use field-normalized citation indicators. Do not
compare impact factors across fields.

\hypertarget{f3-incommensurable-citation-norms}{%
\subsection{F3: Incommensurable citation
norms}\label{f3-incommensurable-citation-norms}}

\textbf{Risk:} Citation rates differ by an order of magnitude across
fields (e.g., mathematics vs.~molecular biology). An integrator bridging
such fields will appear underperforming by the norms of the
high-citation field.

\textbf{Safeguard:} Normalize \(E\) and citation counts by
field-specific baselines. The panel's \(E\) is already a fraction (not
an absolute count), which mitigates this partially.

\hypertarget{f4-misclassification-persistence}{%
\subsection{F4: Misclassification
persistence}\label{f4-misclassification-persistence}}

\textbf{Risk:} A specialist enters the ``Interdisciplinary'' category
and remains there because reclassification is bureaucratically
difficult.

\textbf{Safeguard:} Panel-based reclassification trigger: if
\(\Delta < 0.35\), recommend reclassification to the researcher's
primary field.

\hypertarget{f5-early-career-data-sparsity}{%
\subsection{F5: Early-career data
sparsity}\label{f5-early-career-data-sparsity}}

\textbf{Risk:} Junior researchers have too few publications for reliable
panel values. Nakhoda, Whigham, and Zwanenburg (2023) showed that
Rao-Stirling confidence intervals can span up to 0.6 points for small
samples. Even at mid-career, individual-level estimates carry wider
uncertainty than aggregate measures, making threshold-based
classification inherently noisy.

\textbf{Safeguard:} Minimum publication threshold (e.g., \(n \geq 15\)).
Below this threshold, present panel values with confidence intervals,
flag them as provisional, and defer hard classification unless the full
interval lies inside one profile region.

\hypertarget{f6-gaming-via-strategic-co-authorship}{%
\subsection{F6: Gaming via strategic
co-authorship}\label{f6-gaming-via-strategic-co-authorship}}

\textbf{Risk:} A researcher adds co-authors from distant fields to
inflate \(\Delta\) without genuine integration.

\textbf{Safeguard:} Weight \(\Delta\) by corresponding-author
publications only. Cross-check against \(S\): genuine integration
produces high coherence, while strategic co-authorship does not.

\hypertarget{discussion}{%
\section{Discussion}\label{discussion}}

The protocol proposed here is designed to complement expert judgment,
not replace it. The panel provides structured evidence about the type of
interdisciplinarity, and the committee composition rule ensures that the
right expertise is present. But the final evaluation decision remains
with the committee.

Several limitations should be noted. First, the classification
thresholds are illustrative and require calibration against empirical
distributions of panel values across disciplines and career stages.
Second, the protocol assumes access to citation data, which may not be
available for all researchers (particularly in the humanities). Third,
the mock profiles use simplified data; real researcher portfolios are
messier.

Rafols (2019) argued that science indicators should be contextualized,
multidimensional, and subject to stakeholder validation. The protocol
follows this principle: it is multidimensional (three components),
context-dependent (committee composition adapts to the researcher's
profile), and transparent (thresholds and rules are explicit and
auditable). The panel should remain vector-valued evidence (\(\Delta\),
\(S\), \(E\)), not a single composite score.

The fundamental insight is that ``Interdisciplinary'' is not a single
category. It encompasses integrators, polymaths, and misclassified
specialists --- researchers with qualitatively different profiles that
require qualitatively different evaluation approaches. The panel
provides the resolution to make these distinctions.

\hypertarget{conclusions}{%
\section{Conclusions}\label{conclusions}}

We have proposed an evaluation protocol for national agencies assessing
researchers classified as ``Interdisciplinary.'' The three-component
indicator panel (\(\Delta\), \(S\), \(E\)) provides the structural
characterization needed to distinguish integrators from polymaths and
specialists, compose appropriate evaluation committees, and guard
against six identified failure modes. The protocol is transparent,
auditable, and designed to be adapted to specific national contexts.
Empirical calibration of the classification thresholds against real
researcher distributions is the natural next step.

\hypertarget{references}{%
\section{References}\label{references}}

\begin{itemize}
\tightlist
\item
  Nakhoda, M., Whigham, P., and Zwanenburg, S. (2023). Quantifying and
  addressing uncertainty in the measurement of interdisciplinarity.
  \emph{Scientometrics}, 128:6107--6127.
\item
  Rafols, I. (2019). S\&T indicators in the wild: contextualization and
  participation for responsible metrics. \emph{Research Evaluation},
  28(1):7--22.
\item
  Rivero, A. (2026). Measuring interdisciplinarity: A multi-component
  indicator panel for research evaluation. Manuscript.
\item
  Xiang, S., Romero, D. M., and Teplitskiy, M. (2025). Evaluating
  interdisciplinary research: Disparate outcomes for topic and knowledge
  base. \emph{Proceedings of the National Academy of Sciences},
  122(16):e2409752122.
\end{itemize}

\end{document}
